\documentclass[11pt]{article}

\usepackage[margin=0.85in]{geometry}
\usepackage[T1]{fontenc}
\usepackage[utf8]{inputenc}
\usepackage{lmodern}
\usepackage{amsmath,amssymb,mathtools}
\usepackage{enumitem}
\usepackage{array,booktabs,tabularx}
\usepackage{xcolor}
\usepackage{hyperref}
\usepackage{tcolorbox}
\tcbuselibrary{skins,breakable}
\usepackage{titlesec}

\hypersetup{
    colorlinks=true,
    linkcolor=blue!60!black,
    urlcolor=blue!60!black
}

\definecolor{myblue}{RGB}{20,55,90}
\definecolor{mygray}{RGB}{245,245,245}
\definecolor{mygreen}{RGB}{20,110,70}
\definecolor{myorange}{RGB}{200,120,20}

\titleformat{\section}{\large\bfseries\color{myblue}}{\thesection.}{0.5em}{}
\titleformat{\subsection}{\normalsize\bfseries\color{myblue}}{\thesubsection.}{0.5em}{}

\newtcolorbox{goalbox}{
    colback=blue!3,
    colframe=myblue,
    boxrule=0.8pt,
    arc=2pt
}

\newtcolorbox{warningbox}{
    colback=red!2,
    colframe=red!55!black,
    boxrule=0.8pt,
    arc=2pt
}

\newtcolorbox{notebox}{
    colback=yellow!5,
    colframe=myorange,
    boxrule=0.8pt,
    arc=2pt
}

\setlist[itemize]{nosep,leftmargin=1.5em}
\setlist[enumerate]{nosep,leftmargin=1.5em}

\begin{document}

\begin{center}
    {\LARGE \textbf{Model Lock-In Sheet}}\\[0.4em]
    {\large ECO 317 -- Assignment \#2}\\[0.4em]
    \rule{0.9\linewidth}{0.5pt}
\end{center}

\begin{goalbox}
\textbf{Purpose.} Lock every equation, variable name, timing convention, and solver output \emph{before} any code is written. This sheet is the single source of truth that every teammate and every AI prompt should reference.
\end{goalbox}

\section{How to use this sheet}

\begin{enumerate}
    \item Open the course notes alongside this document.
    \item Confirm that each equation below matches the accessible lecture PDFs and the assignment brief.
    \item Any entry explicitly labeled ``not numerically specified in accessible notes'' should be treated as a coding calibration choice, not as a lecture fact.
    \item During coding and demo prep, reference this sheet instead of re-deriving from memory.
    \item If any equation changes after review, update this sheet first, then update the code.
\end{enumerate}

\section{Shared requirements across all models}

These apply identically to Models 1, 2, and 3.

\begin{tabularx}{\linewidth}{@{}p{2.6cm}X@{}}
\toprule
\textbf{Requirement} & \textbf{Detail} \\
\midrule
Horizon & Infinite \\
Solution method & Value Function Iteration (VFI) \\
Shock process & Two-state Markov chain per model \\
Grid density & Fine grid, approximately 200--400 points \\
Convergence & $\max|V_{n+1}-V_n|<\varepsilon$, e.g.\ $\varepsilon=10^{-6}$ \\
Simulation & At least 100 periods \\
Moments & Mean, variance, lag-1 autocorrelation, correlation with income/output \\
Forecasting & Exact non-linear forecast from solved policy functions; user-selected future shocks \\
Dashboard & Streamlit, sidebar sliders, dark/navy styling, white chart cards, Garamond-style font \\
Summaries & Deterministic f-string text generated from live simulation moments \\
\bottomrule
\end{tabularx}

\bigskip

\section{Model 1: Stochastic consumption-savings}

\subsection{Purpose}
Household chooses how much to save in a risk-free asset when income randomly switches between a bad and a good state.

\subsection{Bellman equation}
Assignment-required infinite-horizon Bellman translation of the lecture's two-period expected-utility problem:
\[
V(a_t,y_t) = \max_{a_{t+1}} \left\{ u(c_t) + \beta \sum_{y_{t+1}} P(y_{t+1}|y_t)\,V(a_{t+1},y_{t+1}) \right\}
\]

\begin{notebox}
Accessible lecture notes present the same economics in two-period bond notation:
\[
\max_{C_t,C_{t+1}} U(C_t) + \beta E_t U(C_{t+1})
\]
subject to
\[
C_t + B_t = Y_t, \qquad C_{t+1} = E(Y_{t+1}) + (1+r_t)B_t.
\]
Use the Bellman form above for coding, but keep the lecture notation in mind when naming variables.
\end{notebox}

\subsection{Budget constraint}
\[
c_t + a_{t+1} = y_t + (1+r_t)\,a_t, \qquad c_t > 0
\]

\begin{notebox}
Lecture notes write saving through a non-state-contingent bond $B_t$ with gross return $(1+r_t)$. In the recursive implementation, the asset state $a_t$ is the same savings vehicle written in grid form.
\end{notebox}

\subsection{Utility function}
CES / CRRA form stated in the assignment:
\[
u(c) =
\begin{cases}
\dfrac{c^{1-\sigma}-1}{1-\sigma}, & \sigma \neq 1 \\[0.5em]
\ln(c), & \sigma = 1
\end{cases}
\]

\subsection{State, control, and shock identification}

\begin{tabularx}{\linewidth}{@{}lX@{}}
\toprule
\textbf{Field} & \textbf{Value} \\
\midrule
State variables & Current assets $a$ and current income shock $y$ \\
Control variable & Next-period assets $a_{t+1}$ \\
Residual variable & Consumption $c_t$ (determined from the budget constraint) \\
Shock & Exogenous income $y \in \{y_L, y_H\}$ \\
Shock interpretation & Low versus high endowment \\
\bottomrule
\end{tabularx}

\subsection{Transition matrix}

\[
P = \begin{pmatrix} \pi_{LL} & \pi_{LH} \\ \pi_{HL} & \pi_{HH} \end{pmatrix}
\qquad\text{where row } i \text{ sums to } 1.
\]

Assignment-required two-state Markov discretization of the lecture's low/high income uncertainty:
\[
P = \begin{pmatrix} 0.9 & 0.1 \\ 0.1 & 0.9 \end{pmatrix}
\]

\begin{notebox}
Lecture notes confirm a two-state uncertainty environment with income outcomes $Y_L$ and $Y_H$ and probabilities $p_L$ and $p_H = 1-p_L$. They do \emph{not} provide a recursive 2-by-2 Markov matrix, so the matrix above is the project's implementation discretization of that lecture setup and can be adjusted by sliders.
\end{notebox}

\subsection{Timing convention}
\begin{itemize}
    \item Period $t$: agent enters with assets $a_t$ and observes income $y_t$.
    \item Agent chooses $a_{t+1}$.
    \item Consumption is determined residually: $c_t = (1+r_t)a_t + y_t - a_{t+1}$.
\end{itemize}

\subsection{Feasibility conditions}
\begin{itemize}
    \item $c_t > 0$ always.
    \item $a_{t+1} \in [a_{\min}, a_{\max}]$.
    \item If a borrowing constraint exists: $a_{t+1} \geq a_{\min}$.
\end{itemize}

\subsection{Parameters}

\begin{tabularx}{\linewidth}{@{}lXl@{}}
\toprule
\textbf{Symbol} & \textbf{Meaning} & \textbf{Coding default / note status} \\
\midrule
$\beta$ & Discount factor & 0.95 \\
$\sigma$ & Risk aversion (CRRA coefficient) & 2.0 \\
$r$ & Risk-free interest rate & 0.03 \\
$y_L$ & Low income state & Implementation calibration chosen to represent the lecture's $Y_L$ state \\
$y_H$ & High income state & Implementation calibration chosen to represent the lecture's $Y_H$ state \\
$a_{\min}$ & Lower bound of asset grid & 0.0 \\
$a_{\max}$ & Upper bound of asset grid & 20.0 \\
$n_a$ & Number of asset grid points & 200--400 \\
\bottomrule
\end{tabularx}

\begin{notebox}
\textbf{Calibration policy for Model 1.} Treat $Y_L$ and $Y_H$ as lecture-confirmed symbols, the infinite-horizon Bellman form as an assignment-required translation, and the numerical choices for $(y_L,y_H,P)$ as implementation calibrations unless another lecture file provides explicit numbers.
\end{notebox}

\subsection{Solver must return}
\begin{itemize}
    \item Value function $V(a,y)$: array of shape \texttt{(n\_a, 2)}.
    \item Policy indices: array of shape \texttt{(n\_a, 2)}.
    \item Savings policy $a_{t+1}(a_t,y_t)$: array of shape \texttt{(n\_a, 2)}.
    \item Consumption policy $c(a,y)$: array of shape \texttt{(n\_a, 2)}.
    \item Convergence diagnostics: iterations, final error, converged flag.
    \item Grid: the asset grid array.
\end{itemize}

\subsection{Simulate, forecast, and plot}
\begin{itemize}
    \item Simulate: $a_t$, $c_t$, $y_t$ for $T \geq 100$ periods.
    \item Forecast: user picks future shocks; apply policy to generate deterministic path.
    \item Plot: consumption path, asset path, income path, policy functions $a_{t+1}(a_t,y_t)$ and $c(a_t,y_t)$.
    \item Moments: mean, variance of $c_t$ and $a_t$; autocorrelation of $c_t$; correlation of $c_t$ with $y_t$.
\end{itemize}

\subsection{Course-note fidelity check}

\begin{warningbox}
Resolved from the accessible lecture notes:
\begin{itemize}
    \item The notes use subscripts such as $C_t$, $Y_t$, and $B_t$, plus the gross return $(1+r_t)$.
    \item The notes describe an endowment/income household that saves through a non-state-contingent bond.
    \item The uncertainty discussion gives two income outcomes $Y_L$ and $Y_H$ and expected utility $E_t U(C_{t+1})$.
    \item The accessible notes do not state a borrowing constraint or numerical calibration for $y_L$, $y_H$, or the transition matrix.
\end{itemize}
\end{warningbox}

\bigskip
\section{Model 2: Stochastic Robinson Crusoe economy}

\subsection{Purpose}
Household accumulates productive capital instead of a risk-free asset. Output depends on capital and a stochastic TFP shock.

\subsection{Bellman equation}
Assignment-required infinite-horizon Bellman translation of the lecture's Robinson Crusoe timing:
\[
V(k_t,z_t) = \max_{k_{t+1}} \left\{ u(c_t) + \beta \sum_{z_{t+1}} P(z_{t+1}|z_t)\,V(k_{t+1},z_{t+1}) \right\}
\]

\begin{warningbox}
\textbf{Locked-in timing.} The lecture notes explicitly say the consumer moves resources forward by choosing $K_{t+1}$ and that tomorrow's output is defined by
\[
Y_{t+1} = F(K_{t+1}).
\]
The Bellman equation above is the assignment-required stochastic infinite-horizon translation of that lecture structure.
\end{warningbox}

\subsection{Resource constraint}
\[
c_t + k_{t+1} = y_t + (1-\delta)\,k_t, \qquad y_{t+1} = z_{t+1}F(k_{t+1}), \qquad c_t > 0
\]

\begin{notebox}
Lecture notes write the two-period resource constraints as
\[
C_t + K_{t+1} = Y_t + (1-\delta)K_t, \qquad C_{t+1} = F(K_{t+1}) + (1-\delta)K_{t+1}.
\]
With the assignment's TFP shock added, the coding convention is: current output is determined by installed capital today, while the choice of $k_{t+1}$ determines tomorrow's output through $y_{t+1} = z_{t+1}F(k_{t+1})$.
\end{notebox}

\subsection{Production function}
Cobb-Douglas (stated in the assignment context):
\[
F(k) = A\,k^\alpha
\]

\subsection{Utility function}
Same CRRA form as Model 1:
\[
u(c) =
\begin{cases}
\dfrac{c^{1-\sigma}-1}{1-\sigma}, & \sigma \neq 1 \\[0.5em]
\ln(c), & \sigma = 1
\end{cases}
\]

\begin{notebox}
The accessible lecture notes impose increasing and concave household utility and use the same consumption theory structure as Model 1, but they do not pin down a new numeric calibration here. Keep the CRRA/log specification above as the coding parameterization unless another lecture note overrides it.
\end{notebox}

\subsection{State, control, and shock identification}

\begin{tabularx}{\linewidth}{@{}lX@{}}
\toprule
\textbf{Field} & \textbf{Value} \\
\midrule
State variables & Current capital $k_t$ and current TFP shock $z_t$ \\
Control variable & Next-period capital $k_{t+1}$ \\
Residual variable & Consumption $c_t$ (from the resource constraint) \\
Derived variables & Output $Y_t = z_tF(k_t)$, investment $I_t = k_{t+1} - (1-\delta)k_t$ \\
Shock & TFP $z \in \{z_L, z_H\}$ \\
Shock interpretation & Bad versus good productivity \\
\bottomrule
\end{tabularx}

\begin{notebox}
Lecture timing is: current output comes from already-installed capital, while the household's choice of $k_{t+1}$ determines next-period output. In the stochastic project version, store current output as $y_t = z_t F(k_t)$ after the state is realized, and carry forward the rule $y_{t+1} = z_{t+1}F(k_{t+1})$.
\end{notebox}

\subsection{Transition matrix}

\[
P^z = \begin{pmatrix} \pi_{LL}^z & \pi_{LH}^z \\ \pi_{HL}^z & \pi_{HH}^z \end{pmatrix}
\]

Assignment-required two-state discretization of the lecture's persistent TFP process:
\[
P^z = \begin{pmatrix} 0.9 & 0.1 \\ 0.1 & 0.9 \end{pmatrix}
\]

\begin{notebox}
The firm notes define TFP as an exogenous AR(1)-style process, $\ln A_t = (1-\rho)A^* + \rho \ln A_{t-1} + \varepsilon_t^a$. Because the assignment requires a two-state Markov shock, $(z_L,z_H,P^z)$ should be interpreted as a discrete approximation of that lecture process unless another note file gives a discrete version directly.
\end{notebox}

\subsection{Timing convention}

\begin{warningbox}
\textbf{Lecture-consistent timing:}
\begin{itemize}
    \item Period $t$: the household enters with predetermined capital $k_t$ and observes current TFP $z_t$.
    \item Current output is generated by installed capital; in the stochastic project version this is stored as $y_t = z_t F(k_t)$.
    \item The household chooses next-period capital $k_{t+1}$.
    \item That choice determines tomorrow's productive capital, so next-period output is $y_{t+1} = z_{t+1}F(k_{t+1})$.
\end{itemize}
\end{warningbox}

\subsection{Feasibility conditions}
\begin{itemize}
    \item $c_t > 0$ always.
    \item $k_{t+1} \in [k_{\min}, k_{\max}]$.
    \item $k_{t+1} \geq 0$ (no negative capital).
\end{itemize}

\subsection{Parameters}

\begin{tabularx}{\linewidth}{@{}lXl@{}}
\toprule
\textbf{Symbol} & \textbf{Meaning} & \textbf{Coding default / note status} \\
\midrule
$\beta$ & Discount factor & 0.95 \\
$\sigma$ & Risk aversion & 2.0 \\
$\alpha$ & Capital share in production & 0.36 \\
$\delta$ & Depreciation rate & 0.10 \\
$A$ & TFP scale (if used) & 1.0 \\
$z_L$ & Low TFP state & Implementation calibration for the low state of the discretized TFP process \\
$z_H$ & High TFP state & Implementation calibration for the high state of the discretized TFP process \\
$k_{\min}$ & Lower bound of capital grid & 0.01 \\
$k_{\max}$ & Upper bound of capital grid & Implementation choice; choose as a multiple of deterministic steady-state capital so the policy does not bind at the upper edge \\
$n_k$ & Number of capital grid points & 200--400 \\
\bottomrule
\end{tabularx}

\begin{notebox}
\textbf{Calibration policy for Model 2.} Treat the lecture AR(1) for TFP and the timing $Y_{t+1}=F(K_{t+1})$ as lecture facts. Treat the two-state objects $(z_L,z_H,P^z)$ and the grid cap $k_{\max}$ as implementation choices used to make the assignment's VFI problem numerically tractable.
\end{notebox}

\subsection{Solver must return}
\begin{itemize}
    \item Value function $V(k,z)$: array of shape \texttt{(n\_k, 2)}.
    \item Policy indices: array of shape \texttt{(n\_k, 2)}.
    \item Capital policy $k_{t+1}(k_t,z_t)$: array of shape \texttt{(n\_k, 2)}.
    \item Consumption policy $c(k_t,z_t)$: array of shape \texttt{(n\_k, 2)}.
    \item Convergence diagnostics.
    \item Grid: the capital grid array.
\end{itemize}

\subsection{Simulate, forecast, and plot}
\begin{itemize}
    \item Simulate: $k_t$, $c_t$, $Y_t$, $I_t$ for $T \geq 100$ periods.
    \item Forecast: user-selected future TFP shocks; apply policy to generate deterministic path.
    \item Plot: capital path, output path, consumption path, investment path, policy functions $k_{t+1}(k_t,z_t)$ and $c(k_t,z_t)$.
    \item Moments: mean and variance of $Y_t$ and $c_t$; autocorrelation of $k_t$ or $Y_t$; correlation of $c_t$ with $Y_t$.
\end{itemize}

\subsection{Course-note fidelity check}

\begin{warningbox}
Resolved from the accessible lecture notes:
\begin{itemize}
    \item The notes explicitly use $Y_{t+1} = F(K_{t+1})$, so the choice of $K_{t+1}$ governs tomorrow's output.
    \item The lecture resource constraints are $C_t + K_{t+1} = Y_t + (1-\delta)K_t$ and $C_{t+1} = F(K_{t+1}) + (1-\delta)K_{t+1}$.
    \item Depreciation enters as undepreciated capital $(1-\delta)K_t$.
    \item The accessible notes define a persistent TFP process, while the assignment is what forces a two-state Markov discretization for implementation.
    \item The accessible notes do not give numerical values for $z_L$, $z_H$, or a project-specific upper grid bound.
\end{itemize}
\end{warningbox}

\bigskip
\section{Model 3: Endogenous labor supply}

\subsection{Purpose}
Household receives a time endowment ($T=1$) and chooses how much labor to supply versus how much leisure to enjoy, with a stochastic wage.

\subsection{Bellman equation}
\[
V(a_t,w_t) = \max_{a_{t+1},\,L_t \in [0,1]} \left\{ u(c_t) + v(1-L_t) + \beta \sum_{w_{t+1}} P(w_{t+1}|w_t)\,V(a_{t+1},w_{t+1}) \right\}
\]

\begin{warningbox}
\textbf{Lecture-confirmed project version.} The accessible lecture notes include labor supply \emph{and} saving through a non-state-contingent bond, so the asset-inclusive version is the lecture-backed default for this project.
\end{warningbox}

\subsection{Budget constraint}
\[
c_t + a_{t+1} = (1+r_t)\,a_t + w_t L_t, \qquad c_t > 0
\]

\begin{notebox}
Accessible lecture notes write the two-period labor problem as
\[
C_t + B_t = w_t L_t - T_t, \qquad C_{t+1} = w_{t+1}L_{t+1} + (1+r_t)B_t - T_{t+1}.
\]
For the project sheet, use the zero-tax recursive version above unless the instructor explicitly wants taxes or profit income carried into the coding version.
\end{notebox}

\subsection{Utility function}
\[
U(c_t,\,1-L_t) = u(c_t) + v(1-L_t),
\qquad
u'(\cdot),\,v'(\cdot) > 0,\quad
u''(\cdot),\,v''(\cdot) < 0
\]

\begin{notebox}
The accessible lecture notes explicitly assume separable utility in consumption and leisure. They do \emph{not} provide a numerical $(\psi,\nu)$ parameterization, so any parametric leisure block used in code should be documented as an implementation choice rather than as a lecture fact.
\end{notebox}

\subsection{State, control, and shock identification}

\begin{tabularx}{\linewidth}{@{}lX@{}}
\toprule
\textbf{Field} & \textbf{Lecture-based project value} \\
\midrule
State variables & Current assets $a_t$ and current wage shock $w_t$ \\
Control variables & Next-period assets $a_{t+1}$ and labor $L_t$ \\
Residual & Consumption $c_t$ (from the budget constraint) \\
Derived & Leisure $\ell_t = 1 - L_t$ \\
Shock & Wage $w_t \in \{w_L, w_H\}$ \\
Interpretation & Low versus high wage \\
\bottomrule
\end{tabularx}

\begin{notebox}
Accessible lecture notes include saving/assets through the bond $B_t$, so the asset-inclusive version is the one locked in here.
\end{notebox}

\subsection{Transition matrix}

\[
P^w = \begin{pmatrix} \pi_{LL}^w & \pi_{LH}^w \\ \pi_{HL}^w & \pi_{HH}^w \end{pmatrix}
\]

Assignment-required two-state wage-shock discretization:
\[
P^w = \begin{pmatrix} 0.9 & 0.1 \\ 0.1 & 0.9 \end{pmatrix}
\]

\begin{notebox}
The accessible notes motivate wage-driven labor income but do not provide a discrete wage-shock calibration. Therefore $(w_L,w_H,P^w)$ should be read as implementation choices used to satisfy the assignment's two-state Markov requirement.
\end{notebox}

\subsection{Timing convention}
\begin{itemize}
    \item Period $t$: the household enters with assets $a_t$ and observes wage $w_t$.
    \item Agent chooses labor $L_t$ and next-period assets $a_{t+1}$.
    \item Consumption is determined residually from the budget constraint.
    \item Time endowment is $T=1$, so leisure is $1-L_t$.
\end{itemize}

\begin{notebox}
Accessible lecture notes also assume an interior solution in which leisure stays strictly between zero and one.
\end{notebox}

\subsection{Feasibility conditions}
\begin{itemize}
    \item $c_t > 0$.
    \item $0 \leq L \leq 1$ (labor bounded by time endowment).
    \item $1 - L > 0$ if utility requires strictly positive leisure.
    \item $a_{t+1} \in [a_{\min}, a_{\max}]$.
\end{itemize}

\subsection{Parameters}

\begin{tabularx}{\linewidth}{@{}lXl@{}}
\toprule
\textbf{Symbol} & \textbf{Meaning} & \textbf{Coding default / note status} \\
\midrule
$\beta$ & Discount factor & 0.95 \\
$\sigma$ & Consumption risk aversion & 2.0 \\
$\psi$ & Weight on leisure in utility & Implementation calibration used only if the coding team chooses a parametric leisure block \\
$\nu$ & Leisure curvature parameter & Implementation calibration used only if the coding team chooses a parametric leisure block \\
$r$ & Interest rate (if assets included) & 0.03 \\
$w_L$ & Low wage state & Implementation calibration for the low state of the assignment's wage shock \\
$w_H$ & High wage state & Implementation calibration for the high state of the assignment's wage shock \\
$n_L$ & Number of labor grid points & 50--100 \\
$n_a$ & Asset grid points (if applicable) & 200--400 \\
\bottomrule
\end{tabularx}

\begin{notebox}
\textbf{Calibration policy for Model 3.} Treat the asset-inclusive budget structure and separable utility as lecture-confirmed. Treat any numerical leisure parameters and the discrete wage-shock values as implementation calibrations unless a later lecture note provides explicit formulas or numbers.
\end{notebox}

\subsection{Solver must return}
\begin{itemize}
    \item Value function $V(a,w)$: shape \texttt{(n\_a, 2)}.
    \item Savings policy $a_{t+1}(a_t,w_t)$.
    \item Labor policy $L(a,w)$.
    \item Consumption policy $c(a,w)$.
    \item Convergence diagnostics.
\end{itemize}

\subsection{Simulate, forecast, and plot}
\begin{itemize}
    \item Simulate: $L_t$, $c_t$, $w_t$, and $a_t$ if applicable, for $T \geq 100$ periods.
    \item Forecast: user-selected future wage shocks; apply policy.
    \item Plot: labor path, consumption path, wage path, assets path (if applicable), static intuition graphs (e.g.\ labor vs.\ wage, labor vs.\ $\psi$).
    \item Moments: mean and variance of $L_t$ and $c_t$; autocorrelation; correlation of $c_t$ with wage income $w_t L_t$.
\end{itemize}

\subsection{Course-note fidelity check}

\begin{warningbox}
Resolved from the accessible lecture notes:
\begin{itemize}
    \item The lecture version includes saving/assets through a bond $B_t$ in addition to labor choice.
    \item Utility is separable in consumption and leisure: $U(c_t,1-L_t)=u(c_t)+v(1-L_t)$.
    \item Leisure enters as $T-L_t$ with $T=1$.
    \item The lecture's two-period constraints use wage income and can include lump-sum taxes; the project sheet uses the corresponding zero-tax recursive form.
    \item The accessible notes do not provide numerical values for $\psi$, $\nu$, $w_L$, $w_H$, or wage-transition probabilities, so those remain documented calibrations.
\end{itemize}
\end{warningbox}

\bigskip
\section{Cross-model coding conventions}

Agree on these variable names before writing any solver, simulation, or dashboard code.

\begin{tabularx}{\linewidth}{@{}llX@{}}
\toprule
\textbf{Code name} & \textbf{Math} & \textbf{Meaning} \\
\midrule
\texttt{beta} & $\beta$ & Discount factor \\
\texttt{sigma} & $\sigma$ & CRRA risk aversion \\
\texttt{r} & $r$ & Risk-free interest rate \\
\texttt{alpha} & $\alpha$ & Capital share (Model 2) \\
\texttt{delta} & $\delta$ & Depreciation rate (Model 2) \\
\texttt{psi} & $\psi$ & Leisure weight (Model 3) \\
\texttt{nu} & $\nu$ & Leisure curvature (Model 3) \\
\midrule
\texttt{a\_grid} & $\{a_1,\ldots,a_{n_a}\}$ & Asset grid \\
\texttt{k\_grid} & $\{k_1,\ldots,k_{n_k}\}$ & Capital grid \\
\texttt{labor\_grid} & $\{L_1,\ldots,L_{n_L}\}$ & Labor grid \\
\midrule
\texttt{y\_vals} & $[y_L,\,y_H]$ & Income shock values (Model 1) \\
\texttt{z\_vals} & $[z_L,\,z_H]$ & TFP shock values (Model 2) \\
\texttt{w\_vals} & $[w_L,\,w_H]$ & Wage shock values (Model 3) \\
\texttt{P} & $P$ & $2\times2$ transition matrix \\
\midrule
\texttt{V} & $V$ & Value function array \\
\texttt{policy\_idx} & -- & Index of optimal control on the grid \\
\texttt{policy\_level} & -- & Level value of optimal control \\
\texttt{c\_policy} & $c(\cdot)$ & Consumption policy array \\
\midrule
\texttt{shock\_idx} & -- & Current shock state index (0 or 1) \\
\texttt{T\_sim} & $T$ & Simulation horizon \\
\texttt{T\_fcast} & -- & Forecast horizon \\
\texttt{seed} & -- & Random seed for reproducibility \\
\bottomrule
\end{tabularx}

\bigskip

\textbf{Solver return format (all models):}

Every solver should return a dictionary with these keys:
\begin{itemize}
    \item \texttt{value\_function}
    \item \texttt{policy\_indices}
    \item \texttt{policy\_levels} (savings/capital/labor as appropriate)
    \item \texttt{c\_policy}
    \item \texttt{grid} (the state grid used)
    \item \texttt{diagnostics} (dict with \texttt{iterations}, \texttt{final\_error}, \texttt{converged})
\end{itemize}

\section{Before coding starts: final checklist}

\begin{enumerate}[label=\textbf{\arabic*.}]
    \item All note-dependent placeholders in this sheet have been resolved or explicitly labeled as not numerically specified in accessible notes.
    \item For Model 2, the timing convention centered on $Y_{t+1} = F(K_{t+1})$ is agreed by the group.
    \item For Model 3, the asset-inclusive lecture version is the one the group intends to code unless another instructor note overrides it.
    \item For Model 3, the utility form is confirmed as separable unless a later course note replaces it.
    \item The coding names in the convention table are agreed upon.
    \item The solver return format is agreed upon.
    \item Every teammate can answer: ``What is the state? What is the control? What is the constraint? What is the timing?'' for each model.
    \item This sheet is saved where every teammate can access it during coding and demo prep.
\end{enumerate}

\begin{goalbox}
\textbf{Rule.} If the course notes disagree with anything on this sheet, the course notes win. Update this sheet, then update the code.
\end{goalbox}

\end{document}
